%% ----------------------------------------------------------------------------
%%  §3.8  Il numero e
%% ----------------------------------------------------------------------------
\section{Il numero $e$}
\label{sec:numero-e}

Il teorema sulle successioni monotòne è utile per definire il \emph{numero di Nepero}
$e$ come limite di una particolare successione monotòna e limitata. Infatti,
introduciamo tale numero mediante il limite
\begin{equation}\label{eq:def-e}
  e=\lim_{n\to+\infty}a_{n}, \qquad\text{con}\qquad a_{n}=\Bigl(1+\frac1n\Bigr)^{n}.
\end{equation}
La \eqref{eq:def-e} è la definizione del \emph{numero di Nepero}\index{numero di Nepero} $e$.
Tale definizione è giustificata dal fatto che, come proveremo tra poco, \emph{la
successione $a_{n}=(1+1/n)^{n}$ è strettamente crescente e limitata}; quindi esiste, ed è
un numero reale, il limite per $n\to+\infty$ di $a_{n}$.

\begin{idea}[Due tendenze opposte]
La successione $(1+1/n)^{n}$ combina fra loro due tendenze opposte: la tendenza della
base a un limite uguale a $1$, e la tendenza dell'esponente a un limite infinito. Se
prevalesse la prima, il limite sarebbe $1$ (come per $1^{n}$); se prevalesse la seconda,
sarebbe $+\infty$ (come per $1{,}01^{n}$, \cref{teo:lim-geometrica}). Il risultato è
invece un numero intermedio tra $1$ e $+\infty$, che non è né l'uno né l'altro: è
proprio questo equilibrio delicato a rendere $e$ così importante. La forma $1^{\infty}$ è
una nuova forma indeterminata, come vedremo alla fine del paragrafo.
\end{idea}

\subsection{Monotonia e limitatezza}

Ci serve una versione «stretta» della disuguaglianza di Bernoulli.

\begin{lemma}[Bernoulli in forma stretta]\label{lem:bernoulli-stretta}
Se $x\ge-1$, $x\neq0$ e $n\ge2$ è un numero naturale, allora $(1+x)^{n}>1+nx$.
\end{lemma}

\begin{schema}
\Ipotesi{$x\ge-1$, \ $x\neq0$; \ $n\in\N$, \ $n\ge2$.}
\Tesi{$(1+x)^{n}>1+nx$.}
\end{schema}

\begin{strategia}
Si ripete la dimostrazione per induzione della disuguaglianza di Bernoulli, partendo da
$n=2$, dove la disuguaglianza è stretta perché $x^{2}>0$. Nel passo induttivo si
moltiplica per $1+x$, che deve essere \emph{positivo} perché la disuguaglianza stretta
si conservi: il caso $x=-1$, in cui $1+x=0$, va quindi controllato a parte.
\end{strategia}

\begin{dimostrazione}
Per $n=2$: $(1+x)^{2}=1+2x+x^{2}>1+2x$, perché $x\neq0$. Supponiamo la tesi vera per un
certo $n\ge2$. Se $x=-1$, si ha $(1+x)^{n+1}=0>1-(n+1)=1+(n+1)x$. Se $x>-1$, moltiplicando
l'ipotesi induttiva per $1+x>0$,
\[
  (1+x)^{n+1}>(1+nx)(1+x)=1+(n+1)x+nx^{2}\ge1+(n+1)x .
\]
Per il principio di induzione (a partire da $n=2$) la tesi vale per ogni $n\ge2$.
\end{dimostrazione}

\begin{teorema}[Esistenza del numero $e$]\label{teo:numero-e}\index{numero di Nepero!esistenza}
Le successioni
\begin{equation}\label{eq:an-bn}
  a_{n}=\Bigl(1+\frac1n\Bigr)^{n}, \qquad b_{n}=\Bigl(1+\frac1n\Bigr)^{n+1}
\end{equation}
sono, rispettivamente, strettamente crescente e strettamente decrescente, e risulta
$a_{n}<b_{n}$ per ogni $n$. Di conseguenza entrambe convergono, allo stesso limite $e$, e
\begin{equation}\label{eq:e-stima}
  a_{n}<e<b_{n} \qquad\perogni n\in\N .
\end{equation}
\end{teorema}

\begin{schema}
\Ipotesi{$a_{n}$ e $b_{n}$ definite dalla \eqref{eq:an-bn}, \ $n\in\N$.}
\Tesi{(i) $a_{n}<a_{n+1}$; \ (ii) $b_{n}>b_{n+1}$; \ (iii) $a_{n}<b_{n}$; \
(iv) $a_{n}$ e $b_{n}$ convergono allo stesso limite $e$, e $a_{n}<e<b_{n}$.}
\end{schema}

\begin{strategia}
Il cuore della dimostrazione sono le due monotonie, ed entrambe si ottengono con lo
stesso trucco. Si scrive il rapporto tra due termini consecutivi, si mettono insieme le
potenze con lo stesso esponente, e si ottiene una disuguaglianza della forma
$\bigl(1+x\bigr)^{n}\gtrless1+(\text{qualcosa})$, con $x=-\frac1{n^{2}}$ per $a_{n}$ e
$x=\frac1{n^{2}-1}$ per $b_{n}$: esattamente la forma della disuguaglianza di Bernoulli.
Poi la conclusione è quasi automatica: $a_{n}$ cresce, $b_{n}$ decresce, e $a_{n}<b_{n}$;
quindi ogni $a_{n}$ è minore di ogni $b_{m}$ (una «tenaglia» come quella di Archimede nel
\cref{sec:premessa-limiti}). In particolare $a_{n}$ è limitata superiormente da $b_{1}=4$,
e il teorema sulle successioni monotòne dà la convergenza. Che i due limiti coincidano
segue dal fatto che $b_{n}=a_{n}(1+\frac1n)$ e $1+\frac1n\to1$.
\end{strategia}

\begin{dimostrazione}
(i) \emph{$a_{n}$ è strettamente crescente.} La tesi è
\begin{equation}\label{eq:an-cresce-1}
  a_{n}=\Bigl(1+\frac1n\Bigr)^{n}>a_{n-1}=\Bigl(1+\frac1{n-1}\Bigr)^{n-1} \qquad\perogni n\ge2 .
\end{equation}
Eseguendo la somma delle frazioni, si può scrivere in modo equivalente
\begin{equation}\label{eq:an-cresce-2}
  \Bigl(\frac{n+1}{n}\Bigr)^{n}>\Bigl(\frac{n}{n-1}\Bigr)^{n}\cdot\frac{n-1}{n},
\end{equation}
cioè, dividendo per $\bigl(\frac{n}{n-1}\bigr)^{n}>0$ e raccogliendo le potenze,
\begin{equation}\label{eq:an-cresce-3}
  \Bigl(\frac{n^{2}-1}{n^{2}}\Bigr)^{n}>\frac{n-1}{n}.
\end{equation}
È conveniente isolare come addendo $1$, nel modo seguente:
\begin{equation}\label{eq:an-cresce-4}
  \Bigl(1-\frac1{n^{2}}\Bigr)^{n}>1-\frac1n .
\end{equation}
La \eqref{eq:an-cresce-4} è equivalente alla tesi da provare; e si ottiene dal
\cref{lem:bernoulli-stretta} con $x=-1/n^{2}$, che è maggiore di $-1$ e diverso da zero:
$\bigl(1-\frac1{n^{2}}\bigr)^{n}>1-n\cdot\frac1{n^{2}}=1-\frac1n$.

(ii) \emph{$b_{n}$ è strettamente decrescente.} Procediamo come nella dimostrazione
precedente. La tesi è
\begin{equation}\label{eq:bn-decresce-1}
  b_{n-1}=\Bigl(1+\frac1{n-1}\Bigr)^{n}>b_{n}=\Bigl(1+\frac1n\Bigr)^{n+1} \qquad\perogni n\ge2,
\end{equation}
cioè, equivalentemente,
\begin{equation}\label{eq:bn-decresce-2}
\begin{gathered}
  \Bigl(\frac{n}{n-1}\Bigr)^{n}>\Bigl(\frac{n+1}{n}\Bigr)^{n}\cdot\frac{n+1}{n},\\
  \text{cioè}\qquad\Bigl(\frac{n^{2}}{n^{2}-1}\Bigr)^{n}>\frac{n+1}{n}.
\end{gathered}
\end{equation}
Come nella dimostrazione precedente isoliamo $1$ come addendo:
\begin{equation}\label{eq:bn-decresce-3}
  \Bigl(1+\frac1{n^{2}-1}\Bigr)^{n}>1+\frac1n .
\end{equation}
La disuguaglianza di Bernoulli, con $x=1/(n^{2}-1)$, dà la conclusione:
\begin{equation}\label{eq:bn-decresce-4}
  \Bigl(1+\frac1{n^{2}-1}\Bigr)^{n}\ge1+\frac{n}{n^{2}-1}>1+\frac1n \qquad\perogni n\ge2
\end{equation}
(nell'ultimo passaggio si è utilizzata la disuguaglianza $\frac{n}{n^{2}-1}>\frac1n$, che
è vera perché equivale a $n^{2}>n^{2}-1$).

(iii) Per ogni $n\in\N$ si ha
\begin{equation}\label{eq:bn-maggiore-an}
  b_{n}=\Bigl(1+\frac1n\Bigr)^{n}\cdot\Bigl(1+\frac1n\Bigr)=a_{n}\cdot\Bigl(1+\frac1n\Bigr)>a_{n}.
\end{equation}

(iv) Dalle (i), (ii), (iii) segue che, per ogni $n,m\in\N$, posto $k=\max\{n,m\}$,
\begin{equation}\label{eq:an-minore-bm}
  a_{n}\le a_{k}<b_{k}\le b_{m};
\end{equation}
in particolare $2=a_{1}\le a_{n}<b_{1}=4$ per ogni $n$. La successione $a_{n}$ è dunque
crescente e limitata, e per il \cref{teo:monotone} converge a un numero reale, che
chiamiamo $e$; analogamente $b_{n}$, decrescente e limitata inferiormente da $a_{1}=2$,
converge. Dalla \eqref{eq:bn-maggiore-an} e dal \cref{teo:operazioni-limiti},
$\lim b_{n}=\lim a_{n}\cdot\lim(1+\frac1n)=e\cdot1=e$. Infine, poiché $a_{n}$ è strettamente
crescente, $a_{n}<a_{n+1}\le\sup_{k}a_{k}=e$; e poiché $b_{n}$ è strettamente
decrescente, $e=\inf_{k}b_{k}\le b_{n+1}<b_{n}$ (si veda la \eqref{eq:lim-monotone}).
\end{dimostrazione}

\begin{strategia}[Un'altra via]
Seguendo \cite{giusti}, si può sviluppare la potenza con la formula del binomio, invece di
confrontare due termini consecutivi. Si ottiene una somma di addendi positivi che
crescono con $n$, e il cui numero cresce con $n$: la crescenza diventa evidente. Ogni
addendo è poi minore di un reciproco di fattoriale, e la limitatezza segue dal confronto
con una progressione geometrica. Si perde la tenaglia \eqref{eq:e-stima}, ma si scopre da
dove viene il numero $e$.
\end{strategia}

\begin{dimostrazione}[Seconda dimostrazione della crescenza e della limitatezza di $a_{n}$]
Per la formula del binomio (\cref{teo:binomio}) con $a=1$ e $b=\frac1n$, e per la
\eqref{eq:combinazioni},
\[
  a_{n}=\sum_{k=0}^{n}\binom nk\frac1{n^{k}}
  =\sum_{k=0}^{n}\frac1{k!}\cdot\frac{n(n-1)\cdots(n-k+1)}{n^{k}} .
\]
La seconda frazione è il prodotto dei $k$ fattori $\frac nn,\frac{n-1}n,\dots,\frac{n-k+1}n$.
Isolando gli addendi con $k=0$ e $k=1$, che valgono $1$, si ottiene
\begin{equation}\label{eq:e-binomio}
  a_{n}=1+1+\sum_{k=2}^{n}\frac1{k!}\Bigl(1-\frac1n\Bigr)\Bigl(1-\frac2n\Bigr)\cdots
  \Bigl(1-\frac{k-1}n\Bigr)
\end{equation}
(per $n=1$ la somma è vuota).

\emph{Crescenza.} Passando da $n$ a $n+1$, ogni fattore $1-\frac jn$, con
$1\le j\le k-1\le n-1$, è positivo e aumenta, perché diventa $1-\frac j{n+1}$; quindi
aumenta ogni addendo della somma \eqref{eq:e-binomio}. Inoltre $a_{n+1}$ contiene un
addendo positivo in più, quello con $k=n+1$. Dunque $a_{n}<a_{n+1}$.

\emph{Limitatezza.} Ogni prodotto nella \eqref{eq:e-binomio} è compreso tra $0$ e $1$.
Inoltre per $k\ge1$ si ha $k!\ge2^{k-1}$, perché ciascuno dei $k-1$ fattori $2,3,\dots,k$
è almeno $2$. Per la \cref{prop:geometrica} con $x=\frac12$ si ha allora
\begin{equation}\label{eq:e-binomio-stima}
  a_{n}\le1+\sum_{k=1}^{n}\frac1{k!}\le1+\sum_{k=1}^{n}\frac1{2^{k-1}}
  =1+\frac{1-(1/2)^{n}}{1-1/2}<3 .
\end{equation}
La successione $a_{n}$ è dunque crescente e limitata superiormente da $3$. Per il
\cref{teo:monotone} converge, e il suo limite $e$ verifica $e\le3$.
\end{dimostrazione}

\begin{osservazione}[Che cosa insegna ciascuna via]\label{oss:e-due-vie}
La prima dimostrazione dà di più sul piano quantitativo: la successione decrescente
$b_{n}$ e la tenaglia $a_{n}<e<b_{n}$, cioè un controllo dell'errore. La seconda dice
invece da dove viene $e$. Posto $s_{n}=\sum_{k=0}^{n}\frac1{k!}$, la
\eqref{eq:e-binomio-stima} dice che $a_{n}\le s_{n}<3$; e con gli stessi strumenti si
prova che $s_{n}\to e$. La successione $s_{n}$ è crescente e limitata, quindi converge a
un limite $s$, e da $a_{n}\le s_{n}$ segue $e\le s$ (\cref{cor:confronto-limiti}).
Viceversa, fissato $m$, per $n\ge m$ si possono trascurare nella \eqref{eq:e-binomio}
gli addendi con $k>m$, che sono positivi:
\[
  a_{n}\ge1+1+\sum_{k=2}^{m}\frac1{k!}\Bigl(1-\frac1n\Bigr)\cdots\Bigl(1-\frac{k-1}n\Bigr).
\]
Per $n\to+\infty$, con $m$ fisso, il secondo membro è una somma di un numero \emph{fisso}
di termini, ciascuno dei quali tende a $\frac1{k!}$ (\cref{teo:operazioni-limiti}):
tende quindi a $s_{m}$, e si ottiene $e\ge s_{m}$. Passando al limite per $m\to+\infty$
si ha $e\ge s$. Dunque $e=\lim_{n}s_{n}$ è la somma della serie dei reciproci dei
fattoriali (\cref{es:serie-fattoriali}), ottenuta qui senza la formula di Taylor. Il punto
delicato è l'ordine dei passaggi al limite: prima $n\to+\infty$ con $m$ fisso, poi
$m\to+\infty$. Nella \eqref{eq:e-binomio} non si può passare al limite addendo per
addendo, perché il numero degli addendi cresce con $n$.
\end{osservazione}

\begin{corollario}\label{cor:e-tra-2-e-3}
Risulta $2<e<3$.
\end{corollario}

\begin{schema}
\Ipotesi{la stima \eqref{eq:e-stima}: $a_{n}<e<b_{n}$ per ogni $n\in\N$.}
\Tesi{$2<e<3$.}
\end{schema}

\begin{strategia}
Basta scegliere nella tenaglia \eqref{eq:e-stima} due indici opportuni: $a_{1}=2$ dà la
stima dal basso; per quella dall'alto occorre un $b_{n}$ minore di $3$, e un calcolo
diretto mostra che $n=5$ è sufficiente.
\end{strategia}

\begin{dimostrazione}
Dalla \eqref{eq:e-stima} con $n=1$ si ha $e>a_{1}=2$, e con $n=5$
\begin{equation}\label{eq:b5}
  e<b_{5}=\Bigl(\frac65\Bigr)^{6}=\frac{46656}{15625}=2{,}985984<3. \qedhere
\end{equation}
\end{dimostrazione}

La \eqref{eq:e-stima} è un'ottima «tenaglia» numerica: aumentando $n$ si ottengono stime
sempre più precise. Riportiamo nella \cref{tab:numero-e} alcuni valori approssimati di
$a_{n}$ e $b_{n}$, rappresentati anche nella \cref{fig:numero-e}. Un metodo molto più efficiente per calcolare espressioni decimali
approssimate di $e$ verrà dalla formula di Taylor (\cref{es:calcolo-e}) e dalla serie dei
reciproci dei fattoriali, $e=\sum_{k=0}^{\infty}\frac1{k!}$ (\cref{oss:e-due-vie},
\cref{es:serie-fattoriali} e \eqref{eq:e-serie}); con esso si ottengono facilmente
valori come
\begin{equation}\label{eq:e-decimali}
  e=2{,}71828182845904523536\ldots
\end{equation}

\begin{table}[htb]
\centering
\renewcommand{\arraystretch}{1.3}
\footnotesize\setlength{\tabcolsep}{3pt}
\begin{tabular}{l*{8}{c}}
\toprule
$n$ & $1$ & $10$ & $50$ & $100$ & $500$ & $1000$ & $10^{5}$ & $10^{6}$\\
\midrule
$(1+1/n)^{n}$ & $2$ & $2{,}5937$ & $2{,}6916$ & $2{,}7048$ & $2{,}7156$ & $2{,}7169$ & $2{,}71827$ & $2{,}718280$\\
$(1+1/n)^{n+1}$ & $4$ & $2{,}8531$ & $2{,}7454$ & $2{,}7319$ & $2{,}7210$ & $2{,}7196$ & $2{,}71830$ & $2{,}718283$\\
\bottomrule
\end{tabular}
\caption{Approssimazioni per difetto ($a_{n}$) e per eccesso ($b_{n}$) del numero $e$.
La convergenza è lenta: con $n=10^{6}$ si ottengono appena cinque cifre decimali
esatte.}
\label{tab:numero-e}
\end{table}

\begin{figure}[htb]
\centering
\includegraphics[width=0.85\textwidth]{numero_e}
\caption{Le successioni $a_{n}=(1+1/n)^{n}$ (crescente, in blu) e
$b_{n}=(1+1/n)^{n+1}$ (decrescente, in rame) stringono il numero $e$ (tratteggiato)
come in una tenaglia (\texttt{python/fig\_numero\_e.py}).}
\label{fig:numero-e}
\end{figure}

\subsection{Generalizzazioni}

Nelle applicazioni sono utili anche i limiti seguenti, generalizzazioni della
definizione \eqref{eq:def-e} a successioni di esponenti reali qualsiasi. Ci serve
preliminarmente un lemma, che permette di «comporre» una successione con una successione
di indici.

\begin{lemma}\label{lem:composizione-indici}
Sia $c_{m}$ una successione con $\lim_{m\to+\infty}c_{m}=L\in\R$, e sia $m_{n}$ una
successione di numeri naturali tale che $m_{n}\to+\infty$. Allora $c_{m_{n}}\to L$.
\end{lemma}

\begin{schema}
\Ipotesi{$c_{m}\to L\in\R$ per $m\to+\infty$; \ $m_{n}\in\N$ e $m_{n}\to+\infty$.}
\Tesi{$c_{m_{n}}\to L$ per $n\to+\infty$.}
\end{schema}

\begin{strategia}
Due definizioni di limite vanno incatenate. Fissato $\epsilon$, la convergenza di $c_{m}$
fornisce una soglia $\mu$ per l'indice $m$; la divergenza di $m_{n}$ garantisce che
definitivamente $m_{n}$ supera quella soglia. Si noti che, a differenza delle estratte,
qui gli indici $m_{n}$ possono ripetersi o non essere crescenti.
\end{strategia}

\begin{dimostrazione}
Fissato $\epsilon>0$, esiste $\mu$ tale che $\abs{c_{m}-L}<\epsilon$ per ogni $m>\mu$.
Poiché $m_{n}\to+\infty$, esiste $\nu$ tale che $m_{n}>\mu$ per ogni $n>\nu$ (si applichi
la \eqref{eq:def-piu-inf} con $M=\max\{\mu,1\}$). Per tali $n$ si ha quindi
$\abs{c_{m_{n}}-L}<\epsilon$.
\end{dimostrazione}

\begin{teorema}\label{teo:e-generalizzato}
Per ogni successione di numeri reali $a_{n}$ (con $a_{n}\neq0$ e $1+1/a_{n}>0$) si ha
\begin{align}
  a_{n}\to+\infty \quad&\implica\quad \Bigl(1+\frac1{a_{n}}\Bigr)^{a_{n}}\to e; \label{eq:e-piu-inf}\\
  a_{n}\to-\infty \quad&\implica\quad \Bigl(1+\frac1{a_{n}}\Bigr)^{a_{n}}\to e. \label{eq:e-meno-inf}
\end{align}
\end{teorema}

\begin{schema}
\Ipotesi{$a_{n}\neq0$, \ $1+\frac1{a_{n}}>0$; \ $a_{n}\to+\infty$ oppure
$a_{n}\to-\infty$.}
\Tesi{$\Bigl(1+\dfrac1{a_{n}}\Bigr)^{a_{n}}\to e$.}
\end{schema}

\begin{strategia}
Per la \eqref{eq:e-piu-inf} confrontiamo con il caso intero mediante la parte intera
$m_{n}=\lfloor a_{n}\rfloor$: poiché $m_{n}\le a_{n}<m_{n}+1$, la base $1+\frac1{a_{n}}$ è
compresa tra $1+\frac1{m_{n}+1}$ e $1+\frac1{m_{n}}$, e l'esponente tra $m_{n}$ e
$m_{n}+1$. Prendendo la base più piccola con l'esponente più piccolo, e la base più
grande con l'esponente più grande, otteniamo due carabinieri che sono, a meno di
fattori che tendono a $1$, proprio le successioni $a_{m}$ e $b_{m}$ del teorema
precedente, valutate in $m=m_{n}\to+\infty$. Per la \eqref{eq:e-meno-inf} un cambiamento
di variabile, $b_{n}=-a_{n}-1$, la riconduce alla precedente.
\end{strategia}

\begin{dimostrazione}[Dimostrazione della \eqref{eq:e-piu-inf}]
Poiché $a_{n}\to+\infty$, definitivamente $a_{n}\ge1$; consideriamo solo tali indici.
Indichiamo con $m_{n}=\lfloor a_{n}\rfloor$ la parte intera di $a_{n}$
(\cref{prop:parte-intera}), il più grande intero minore o uguale ad $a_{n}$: risulta
$m_{n}\in\N$, $m_{n}\le a_{n}<m_{n}+1$, e quindi $m_{n}>a_{n}-1\to+\infty$. Si ha
\begin{equation}\label{eq:e-carabinieri}
  \Bigl(1+\frac1{m_{n}+1}\Bigr)^{m_{n}}\le\Bigl(1+\frac1{a_{n}}\Bigr)^{a_{n}}\le\Bigl(1+\frac1{m_{n}}\Bigr)^{m_{n}+1},
\end{equation}
perché la base e l'esponente del primo membro sono minori o uguali rispettivamente della
base e dell'esponente del secondo membro, e lo stesso accade tra secondo e terzo membro:
si usano la monotonia della potenza rispetto all'esponente per basi maggiori di $1$
\eqref{eq:pot5} e la monotonia rispetto alla base per esponenti positivi \eqref{eq:pot3}.
Il risultato segue dal teorema dei carabinieri, perché il primo e il terzo membro tendono
ad $e$. Infatti, per il terzo membro, $\bigl(1+\frac1{m}\bigr)^{m+1}=b_{m}\to e$; per il
primo,
\begin{equation}\label{eq:e-primo-membro}
  \Bigl(1+\frac1{m+1}\Bigr)^{m}=\frac{\bigl(1+\frac1{m+1}\bigr)^{m+1}}{1+\frac1{m+1}}
  =\frac{a_{m+1}}{1+\frac1{m+1}}\to\frac e1=e ;
\end{equation}
in entrambi i casi per $m\to+\infty$, e quindi anche per $m=m_{n}$, per il
\cref{lem:composizione-indici}.
\end{dimostrazione}

\begin{dimostrazione}[Dimostrazione della \eqref{eq:e-meno-inf}]
Poniamo
\begin{equation}\label{eq:e-cambio}
  b_{n}=-a_{n}-1 .
\end{equation}
Dato che $a_{n}\to-\infty$, la successione $b_{n}$ diverge a $+\infty$; in particolare
definitivamente $b_{n}>0$. Inoltre $a_{n}=-(b_{n}+1)$, e quindi
\begin{equation}\label{eq:e-meno-inf-passo}
\begin{split}
  \Bigl(1+\frac1{a_{n}}\Bigr)^{a_{n}}&=\Bigl(1-\frac1{b_{n}+1}\Bigr)^{-(b_{n}+1)}
  =\Bigl(\frac{b_{n}}{b_{n}+1}\Bigr)^{-(b_{n}+1)}\\
  &=\Bigl(\frac{b_{n}+1}{b_{n}}\Bigr)^{b_{n}+1}
  =\Bigl(1+\frac1{b_{n}}\Bigr)^{b_{n}}\cdot\Bigl(1+\frac1{b_{n}}\Bigr).
\end{split}
\end{equation}
Ricordando che $b_{n}\to+\infty$, l'ultimo membro converge a $e\cdot1=e$, a causa della
\eqref{eq:e-piu-inf} e della \eqref{eq:inf5}; ciò prova la \eqref{eq:e-meno-inf}.
\end{dimostrazione}

\begin{osservazione}
Nelle due dimostrazioni, e in alcune delle successive, abbiamo usato le proprietà delle
potenze a esponente reale. La monotonia rispetto all'esponente \eqref{eq:pot5} è stata
dimostrata nel \cref{teo:esponenziale}(ii). Le regole algebriche \eqref{eq:pot1} e la
monotonia rispetto alla base \eqref{eq:pot3}, invece, per esponenti irrazionali sono
state soltanto discusse nel \cref{sec:esponenziale-reale}: la loro dimostrazione
completa verrà data quando avremo a disposizione la continuità della funzione
esponenziale.
\end{osservazione}

\begin{esempio}\label{es:limiti-e}
Dal \cref{teo:e-generalizzato} e dalle operazioni con i limiti seguono molti limiti
utili. Ad esempio:
\begin{itemize}
  \item $\bigl(1+\frac1n\bigr)^{2n}=\bigl[\bigl(1+\frac1n\bigr)^{n}\bigr]^{2}\to e^{2}$;
  \item $\bigl(1-\frac1n\bigr)^{n}\to\frac1e$: infatti, per la \eqref{eq:e-meno-inf} con
  $a_{n}=-n$ (per $n\ge2$, così che $1+1/a_{n}>0$), $\bigl(1-\frac1n\bigr)^{-n}\to e$, e il
  reciproco tende a $1/e$;
  \item $\bigl(\frac{n+1}{n-1}\bigr)^{n}=\frac{(1+1/n)^{n}}{(1-1/n)^{n}}\to\frac{e}{1/e}=e^{2}$
  (anche qui $n\ge2$).
\end{itemize}
\end{esempio}

\subsection{Nuove forme indeterminate}

La forma $1^{\infty}$ è una nuova forma indeterminata, come pure sono indeterminate le
forme $\infty^{0}$ e $0^{0}$. Quindi all'elenco \eqref{eq:forme-indeterminate} vanno
aggiunte le \emph{forme indeterminate}\index{forma indeterminata!esponenziale}
\begin{equation}\label{eq:forme-esponenziali}
  1^{(+\infty)}, \qquad 1^{(-\infty)}, \qquad (+\infty)^{0}, \qquad 0^{0} .
\end{equation}
Che la forma $1^{\infty}$ sia davvero indeterminata lo mostrano tre successioni in cui la
base tende a $1$ e l'esponente a $+\infty$, con limiti diversi:
\begin{gather*}
  \Bigl(1+\frac1n\Bigr)^{n}\to e, \qquad
  \Bigl(1+\frac1n\Bigr)^{n^{2}}\ge1+n^{2}\cdot\frac1n=1+n\to+\infty,\\
  1\le\Bigl(1+\frac1{n^{2}}\Bigr)^{n}=\Bigl[\Bigl(1+\frac1{n^{2}}\Bigr)^{n^{2}}\Bigr]^{1/n}<3^{1/n}\to1 .
\end{gather*}
Nella seconda si è usata la disuguaglianza di Bernoulli, nella terza la stima
$a_{n^{2}}<3$ del \cref{cor:e-tra-2-e-3} e la \cref{prop:radice-n-a}.

\begin{notastorica}[Il numero $e$ da Eulero a oggi]
Nell'\emph{Introductio in analysin infinitorum} (1748) Eulero mostrò che il limite
\eqref{eq:def-e} coincide con la somma della serie
$1+1+\frac1{2!}+\frac1{3!}+\frac1{4!}+\cdots$, che converge molto più rapidamente (la
ritroveremo nel \cref{cap:serie}, formula \eqref{eq:e-serie}), e ne calcolò le prime $23$ cifre decimali; nella
stessa opera introdusse la funzione esponenziale $e^{x}$ e il logaritmo naturale come sua
inversa, nella forma che usiamo ancora oggi. Il fatto che la successione $(1+1/n)^{n}$
converga così lentamente (si veda la \cref{tab:numero-e}) è una curiosità ben nota agli
studenti di calcolo numerico: con i numeri in virgola mobile a doppia precisione, per
$n$ molto grande il calcolo di $1+1/n$ perde tutte le cifre significative di $1/n$, e il
risultato «converge» a $1$ invece che ad $e$ (si veda il programma qui sotto).
\end{notastorica}

\begin{python}[Il numero $e$ con la tenaglia, e un'insidia numerica]
import math

for n in [10, 1000, 10**5, 10**7]:
    a = (1 + 1/n) ** n
    b = (1 + 1/n) ** (n + 1)
    print(f"n = {n:>8d}:  {a:.10f} < e < {b:.10f}"
          f"   (ampiezza {b - a:.1e})")
print(f"math.e        = {math.e:.10f}")

# Insidia della virgola mobile: per n enorme 1 + 1/n si arrotonda.
# Confrontiamo il calcolo "ingenuo" con quello basato su log1p.
for k in [9, 12, 15, 16]:
    n = 10**k
    ingenuo = (1 + 1/n) ** n
    accurato = math.exp(n * math.log1p(1/n))
    print(f"n = 1e{k}:  ingenuo = {ingenuo:.10f}"
          f"   con log1p = {accurato:.10f}")
\end{python}
\begin{risultato}
n =       10:  2.5937424601 < e < 2.8531167061   (ampiezza 2.6e-01)
n =     1000:  2.7169239322 < e < 2.7196408562   (ampiezza 2.7e-03)
n =   100000:  2.7182682372 < e < 2.7182954199   (ampiezza 2.7e-05)
n = 10000000:  2.7182816941 < e < 2.7182819660   (ampiezza 2.7e-07)
math.e        = 2.7182818285
n = 1e9:  ingenuo = 2.7182820520   con log1p = 2.7182818271
n = 1e12:  ingenuo = 2.7185234960   con log1p = 2.7182818285
n = 1e15:  ingenuo = 3.0350352065   con log1p = 2.7182818285
n = 1e16:  ingenuo = 1.0000000000   con log1p = 2.7182818285
\end{risultato}

Si noti che già per $n=10^{9}$ il valore calcolato «ingenuamente» supera $e$, in
contraddizione con la \eqref{eq:e-stima}: l'errore di arrotondamento nel calcolo di
$1+1/n$, amplificato dall'elevamento alla potenza $n$, ha superato l'errore di
troncamento. Per $n=10^{15}$ il risultato è addirittura maggiore di $3$, e per $n=10^{16}$
vale $1$, perché $1+10^{-16}$ viene arrotondato a $1$. La matematica dice che la
successione converge ad $e$; il calcolatore, se usato senza cautela, dice altro. Il
rimedio, in questo caso, è riscrivere la potenza come $e^{\,n\log(1+1/n)}$ e calcolare
$\log(1+x)$ con una funzione apposita (\texttt{math.log1p}), che non perde cifre per $x$
piccolo: la terza colonna mostra che così i valori restano corretti.
